% !TEX root = tarasov-qasp-zh.tex

\chapter[你知道欧姆定律吗？]{你知道欧姆定律吗？\label{ch-27}}

\begin{dialogue}

\tchr 你知道欧姆定律吗？

\stda 当然知道。我想谁都知道欧姆定律。这大概是整个物理课程里最简单的问题了。

\tchr 那我们来看看。\figr{fig-104}~\drkgry{(a)} 所示是电路中的一段（portion of a circuit）。其中 $\Ea$ 是电动势（electromotive force，简称 emf），方向向右；$R_{1}$、$R_{2}$ 是电阻；$r$ 是电动势源（seat of the electromotive force，即电源）的内阻（internal resistance）；$\varphi_{A}$、$\varphi_{B}$ 是所给这段电路两端的电势。电流自左向右流动。求此电流的值 $I$。
\begin{figure}%
\centering
\includegraphics[width=0.6\textwidth]{figs/ch-27/fig-104a}
\caption[求电路中的电流。]{求电路中的电流。\label{fig-104}}
\end{figure}

\stda 可是您给出的是一段断开的电路呀！

\tchr 我让你考虑的是某个更大电路中的一段。电路的其余部分你们一无所知，也无需知道，因为这一段两端的电势已经给定。

\stda 以前我们只处理过闭合电路（closed circuit）。对闭合电路，欧姆定律可以写成如下形式：
\begin{equation}%
I = \frac{\Ea}{R + r}
\label{eq-164}
\end{equation}

\tchr 你错了。你们也处理过一段电路（element of a circuit）的问题。按照欧姆定律，一段电路中的电流等于电压与电阻之比。

\stda 可这算一段电路吗？

\tchr 当然算。\figr{fig-104}~\drkgry{(b)} 画的就是这样一段。对这一段，欧姆定律可以写成
\begin{equation}%
I = \frac{\varphi_{A} - \varphi_{B}}{R}
\label{eq-165}
\end{equation}
对于这段电路两端的电势差 $(\varphi_{A} - \varphi_{B})$，你们以前用的是更简单的名称——“电压（voltage）”，用字母 $V$ 表示。

\stda 无论如何，\figr{fig-104}~\drkgry{(a)} 所示的这种一段电路，我们是没有处理过的。

\tchr 这么说，欧姆定律你只在两种特殊情形下会用：闭合电路，以及不含电动势的最简单的一段；一般情形下的欧姆定律你并不知道。我们一起来把它弄清楚。

\figr{fig-105}~\drkgry{(a)} 画出了电势沿这段电路的变化。电流自左向右流动，因此电势从 $A$ 到 $C$ 下降。电阻 $R_{1}$ 上的电势降（potential drop）等于 $IR_{1}$。再假定原电池（galvanic cell）的两块极板位于 $C$ 和 $D$。在这两点处发生向上的电势跃变（potential jump），跃变之和即电动势，等于 $\Ea$。在 $C$、$D$ 之间，电势沿电池的内阻下降；电阻 $R_{2}$ 上的电势降等于 $IR_{2}$。这段电路所有电阻上的电势降之和，减去向上的电势跃变，就等于 $V$，也就是所考察的这段电路两端之间的电势差。于是
\begin{equation*}%
I(R_{1}+R_{2}+r) - \Ea=  \varphi_{A} - \varphi_{B}
\end{equation*}
\begin{figure}[h]%
\centering
\includegraphics[width=\textwidth]{figs/ch-27/fig-105a}
\caption[电势沿一段电路的变化。]{电势沿一段电路的变化：(a) 电动势与电流同向（对应图~\ref{fig-104}~\drkgry{(a)}）；(b) 电动势与电流反向（对应图~\ref{fig-104}~\drkgry{(c)}）。\label{fig-105}}
\end{figure}
\pnote{底本\figr{fig-105} 与\figr{fig-106} 的图题均误沿了\figr{fig-104} 的文字 ``Find the value of the current in the circuit''，与两图内容（电势沿一段电路的变化；电压表误差计算的电路）不符，显系录入之误；中文版按图面内容拟写图题，待考原书。}
\pnote{“电势跃变”发生在电池内部：电极与电解质接触处形成厚度仅纳米量级的双电层，电池的非静电力（化学作用）正是在这极薄的层内把电荷送到高电势处。人教版不画电势分布图，而在必修第三册第十二章以“非静电力做功”引入电动势。沿回路把各段电势降与电势跃变依次相加，就得到大学电路课程里的基尔霍夫第二定律（回路电压方程）；\figr{fig-105} 这类图线正是它最直观的形式。}

由此得出电流的表达式，即适用于电路这一段的欧姆定律
\begin{equation}%
I = \frac{\Ea + (\varphi_{A} - \varphi_{B})}{R_{1}+R_{2}+r}
\label{eq-166}
\end{equation}
注意，由这最后一个方程很容易得出你们熟悉的几种特殊情形。对不含电动势的最简单一段，在方程~\eqref{eq-166} 中代入 $\Ea=0$ 和 $r=0$，得
\begin{equation*}%
I = \frac{\varphi_{A} - \varphi_{B}}{R_{1}+R_{2}}
\end{equation*}
它与方程~\eqref{eq-165} 相对应。要得到闭合电路，须把我们这一段的 $A$、$B$ 两端连接起来，这就意味着 $\phi_{A} = \phi_{B}$，于是
\begin{equation*}
I = \frac{\Ea}{(R_{1}+R_{2}+r)}
\end{equation*}
这与方程~\eqref{eq-164} 相对应。
\pnote{式~\eqref{eq-166} 通常称为一段含源电路的欧姆定律。中国教材只讲它的两个特例：初中的欧姆定律 $I=U/R$（人教版九年级第十七章，即式~\eqref{eq-165}，旧称“部分电路的欧姆定律”），以及人教版必修第三册第十二章第 2 节《闭合电路的欧姆定律》（式~\eqref{eq-164}）。符号对照：人教版记电动势为 $E$、电压为 $U$，对应原书的花体 $\Ea$ 与 $V$。一般含源电路的欧姆定律及其符号规则（电动势顺着电流取正、逆着取负，见后文式~\eqref{eq-168}），到大学普通物理的基尔霍夫定律中才系统展开。}

\stda 现在我明白了，我原先其实并不懂欧姆定律。

\tchr 更确切地说，你只是在几种特殊情形下懂它。设想在\figr{fig-104}~\drkgry{(a)} 所示的一段电路中，把一个电压表（voltmeter）接到电池的两端，并设电压表的电阻足够大，它接入电路所引起的畸变可以忽略。电压表将指示什么？

\stda 我知道，接在电池两端的电压表，指示的应当是外电路（external circuit）上的电压降。可是在眼下这种情形，外电路的情况我们一无所知。

\tchr 就我们的目的而言，并不需要知道外电路的情况。电压表若接在 $C$、$D$ 两点，它指示的就是这两点之间的电势差。这一点你明白吧？

\stda 是的，当然明白。

\tchr 现在来看\figr{fig-105}~\drkgry{(a)}。显然，$C$、$D$ 两点间的电势差等于 $(\Ea - Ir)$。用 $V$ 表示电压表的读数，就得到公式
\begin{equation}%
V = \Ea - Ir
\label{eq-167}
\end{equation}
我建议你们就使用这个公式，因为它不需要知道任何外电阻。在处理或多或少比较复杂的电路时，这一点特别有价值。
\pnote{底本此句原作 ``equals $(\Ea = Ir)$''，与紧随其后的式~\eqref{eq-167} 及\figr{fig-105}~\drkgry{(a)} 均不符，显系 ``$-$'' 误植为 ``$=$''，中文版径改为 $(\Ea - Ir)$。本章另有若干文法与标点级小误（如 ``flows form'' 应为 ``from''、``the drop on potential'' 应为 ``in potential''、``write Ohm's in the form'' 脱一 ``law'' 字、式~\eqref{eq-166} 后衍一右括号、陈述句句末误用问号等），均径改不注。}
注意，方程~\eqref{eq-167} 正是一条著名规则的依据：若电路断开、没有电流（$I = 0$），则 $V =\Ea$，电压表的读数与电动势的数值一致。这些你都清楚了吗？

\stda 是的，现在我清楚了。

\tchr 作为检验，我来问你一个问题，考生们常常答不好它。\hlred{闭合电路由 $n$ 个串联的电池组成，每个电池的电动势为 $\Ea$、内阻为 $r$，连接导线的电阻设为零。接在其中一个电池两端的电压表，读数将是多少？照例设电压表中不通过电流。}
\pnote{式~\eqref{eq-167} 即人教版必修第三册第十二章第 2 节的路端电压公式 $U=E-Ir$：路端电压是外电路上的电压，随电流增大而线性减小；电路断开时它等于电动势，这正是用电压表粗测电动势的依据。随后的第 3 节实验再由多组 $(U,I)$ 数据作图定出 $E$ 与 $r$（图线在 $U$ 轴上的截距为 $E$、斜率大小为 $r$），比单次断路测量可靠得多。}

\stda 我按照前面讲解的思路来推理。电压表的读数为 $V = \Ea -Ir$。由这个电路的欧姆定律可以求出电流
\begin{equation*}
I = \frac{n\Ea}{nr} = \frac{\Ea}{r}
\end{equation*}
把它代入第一个式子，得
\begin{equation*}
V = \Ea - \left( \frac{\Ea}{r} \right)r = 0
\end{equation*}
可见这种情形下电压表读数为零。

\tchr 完全正确。只是请记住，这个情形是理想化了的：一方面忽略了连接导线的电阻，另一方面又假定电压表的电阻无穷大。所以，不要试图用实验来检验这个结果。
\pnote{$V=0$ 的由来：$n$ 个相同电池串联，总电动势为 $n\Ea$、总内阻为 $nr$，故 $I=\Ea/r$ 与电池个数无关；而每个电池的路端电压 $\Ea-Ir$ 恰好为零——电动势全部消耗在自身的内阻上。教师“不要用实验检验”的告诫可以定量表述：设连接导线的电阻为 $R_{0}$，则 $I=\dfrac{n\Ea}{nr+R_{0}}$，电压表读数为 $\Ea\cdot\dfrac{R_{0}}{nr+R_{0}}$，是很小却不为零的值；再计入电压表有限的电阻，偏离还要更大些。}

现在来考虑这样的情形：一段电路中电流朝一个方向流动，而电动势沿相反方向作用。\figr{fig-104}~\drkgry{(c)} 画的就是这种情形。请画出电势沿这一段变化的图线。

\stda 电流怎么可能逆着电动势流动呢？

\tchr 你忘了，我们这里有的只是电路的一段。在所考察的这段之外，电路中可能还有别的电动势，在它们的作用下，这一段中的电流就可能逆着给定的电动势流动。

\stda 我明白了。既然电流自左向右流动，从 $A$ 到 $C$ 有等于 $IR_{1}$ 的电势降。既然电动势如今方向相反，$C$、$D$ 两点处的电势跃变就应当是降低电势，而不是升高电势。从 $C$ 到 $D$，电势应下降 $Ir$；从 $D$ 到 $B$，下降 $IR_{2}$。这样就得到了\figr{fig-105}~\drkgry{(b)} 的图线。

\tchr 那么在这种情形下欧姆定律取什么形式？

\stda 它将具有如下形式
\begin{equation}%
I = \frac{ (\varphi_{A} - \varphi_{B}) - \Ea}{R_{1}+R_{2}+r}
\label{eq-168}
\end{equation}

\tchr 正确。那么现在电压表将指示什么？

\stda 从\figr{fig-105}~\drkgry{(b)} 可以看出，这时
\begin{equation}%
V = \Ea + Ir
\end{equation}

\tchr 正是如此。现在来考虑下面的问题。\hlred{在\figr{fig-106} 所示的电路中，$r=\SI{1}{\ohm}$，$R = \SI{10}{\ohm}$，电压表的电阻 $R_{v} = \SI{200}{\ohm}$。计算电压表读数的相对误差——比较基准是设电压表电阻为无穷大、从而不使电路发生畸变时的读数。}

\begin{figure}[!h]%
\centering
\includegraphics[width=0.4\textwidth]{figs/ch-27/fig-106a}
\caption[计算电压表读数的相对误差。]{计算电压表读数相对误差用的电路。\label{fig-106}}
\end{figure}

我们把实际电压表的读数记作 $V$，电阻为无穷大的电压表的读数记作 $V_{\infty}$。于是相对误差为
\begin{equation}
f  = \frac{V_{\infty} - V}{V_{\infty}} = 1 - \frac{V}{V_{\infty}}
\label{eq-170}
\end{equation}
其次，我们考虑到
\begin{equation}
V_{\infty} = \left( \frac{\Ea}{R+r} \right) \,R
\label{eq-171}
\end{equation}
以及
\begin{equation}
V = \frac{\Ea}{r + \left( \dfrac{RR_{V}}{R + R_{V}} \right)} \left( \frac{RR_{V}}{R+R_{V}} \right)
\label{eq-172}
\end{equation}
把方程~\eqref{eq-171} 和~\ref{eq-172} 代入式~\eqref{eq-170}，得
\begin{align*}
f & = 1 - \frac{R_{V} (R + r)}{(R + R_{V})r + RR_{V}} \\
& = 1 - \frac{R_{V} (R + r)}{(r+R)R_{V} + rR} \\
&= 1 - \frac{1}{1+ \left( \dfrac{rR}{(r+R)R_{V}} \right)}
\end{align*}
由于 $R_{V}\gg R$ 且 $R> r$，最后一个式子分母中的那个分数远小于 1。因此，可以利用一个永远值得记住的近似公式
\begin{equation}
(1+\lambda)^{\alpha} \cong 1 + \alpha \lambda
\label{eq-173}
\end{equation}
这个公式在 $\lambda \ll 1$ 时对任何 $\alpha$（整数或分数、正数或负数）都成立。运用近似公式~\eqref{eq-173}，取
\begin{equation*}
\alpha = -1, \quad \lambda = \frac{rR}{(r+R)R_{V}}
\end{equation*}
便得
\begin{equation}
f \cong \frac{rR}{(r+R)R_{V}}
\label{eq-174}
\end{equation}
把给定的数值代入方程~\eqref{eq-174}，求得误差为 $f \approx \frac{1}{220} = 0.0045$。

\stda 这是不是意味着：电压表的电阻比外电阻高出越多，相对误差就越小，我们也就越有理由忽略电压表接入电路所引起的畸变？

\tchr 是的，是这样。只是要记住：$R_{V} \gg R$ 是误差 $f$ 很小的充分条件，但不是必要条件。由方程~\eqref{eq-174} 可见，只要满足条件 $R_{V}\gg r$——即电压表的电阻远大于电源的内阻——误差 $f$ 就很小。这时外电阻甚至可以是无穷大。
\pnote{由式~\eqref{eq-174} 可得更紧的估计：恒有 $\dfrac{rR}{r+R}<r$，故 $f<\dfrac{r}{R_{V}}$——只要 $R_{V}\gg r$，误差必然很小，与外电阻多大无关；教师“外电阻甚至可以无穷大”一句的定量含义即此（$R\to\infty$ 时 $f\to r/R_{V}$）。与这一分析同源的高考常规考点，是“伏安法测电阻”中电流表内接、外接的系统误差选择。另备考：题设 $R_{V}=\SI{200}{\ohm}$ 是当年指针式电压表的量级，现代数字电压表的输入阻抗典型达 \SI{10}{\mega\ohm} 量级，这项误差在实践中已微乎其微；但“电表不理想”的模型仍是高考实验题与竞赛的常见题型（本章最后两道习题即其变式）。}

请试着解下面的问题：
\hlred{在\figr{fig-107}~\drkgry{(a)} 所示的电路中，$E= \SI{6}{\volt}, \, r = 2/3\;\Omega, \, R = \SI{2}{\ohm}$。计算电压表的读数。}

\begin{figure}%
\centering
\includegraphics[width=0.5\textwidth]{figs/ch-27/fig-107a}
\caption[求电压表的读数。]{求电压表的读数。\label{fig-107}}
\end{figure}

\stda 可以认为电压表的电阻是无穷大吗？

\tchr 可以，何况题目里根本没有给出这个电阻。

\stda 可是这样一来，电流还会流过电路中间的那些电阻吗？它多半会直接沿 $A_{1}A_{2}$ 和 $B_{1}B_{2}$ 这两段流过去吧。

\tchr 你错了。在讨论电流之前，我建议你先把图适当化简。由于 $A_{1}A_{2}$ 和 $B_{1}B_{2}$ 这两段没有电阻，所以 $\varphi_{A1} = \varphi_{A2}$，$\varphi_{B1} = \varphi_{B2}$。接下来可以利用这样一条规则：电路中任意两点只要电势相同，就可以把它们合并成一点，而不改变通过各电阻的电流。把这条规则用到我们的情形：让 $A_{1}$ 与 $A_{2}$ 重合、$B_{1}$ 与 $B_{2}$ 重合，就得到\figr{fig-107}~\drkgry{(b)} 所示的电路。它已经很好处理了。因此我直接给出最终答案：电压表的读数是 \SI{4}{\volt}。必要的计算留给你，作为课外作业。
\pnote{底本此句有脱漏：``if in a circuit any two points have the same potential, without changing the currents through the resistors'' 缺主要谓语，句不成句；按文意补足为“电势相同的任意两点可以合并成一点，而不改变通过各电阻的电流”。道理很直白：两点既然等势，即使再用导线把它们连起来，这根导线中也不会有电流，故合并与否不影响电路。“等势点缩点”与“无电阻导线两端等势、可缩为一个节点”是同一条等效变换规则的两种说法，是化简看似复杂的电路（如\figr{fig-107}~\drkgry{(a)}）的基本手段，与串、并联的等效电阻公式（人教版必修第三册第十一章第 4 节）配合使用。}

\end{dialogue}

\section*{习题}
\label{prob-27}
\addcontentsline{toc}{section}{习题}

\begin{enumerate}[resume*=problems]
\item 在\figr{fig-108} 所示电路的一条支路中接入的电流表（ammeter），读数为 \SI{0.5}{\ampere}。求通过电阻 $R_{4}$ 的电流，已知各电阻为：$R_{1} = \SI{2}{\ohm}$、$R_{2}= \SI{4}{\ohm}$、$R_{3}= \SI{1}{\ohm}$、$R_{4}= \SI{2}{\ohm}$、$R_{5}= \SI{1}{\ohm}$。
\begin{figure}[h]%
\centering
\includegraphics[width=0.5\textwidth]{figs/ch-27/fig-108a}
\caption[求通过指定电阻的电流。]{求通过电阻 $R_{4}$ 的电流。\label{fig-108}}
\end{figure}

\item 在\figr{fig-109} 所示的电路中，$\Ea=\SI{4}{\volt}$、$r= \SI{1}{\ohm}$、$R= \SI{2}{\ohm}$。求电流表的读数。
\begin{figure}%
\centering
\includegraphics[width=0.35\textwidth]{figs/ch-27/fig-109a}
\caption[求电流表的读数。]{求电流表的读数。\label{fig-109}}
\end{figure}

\item 检流计（galvanometer）的电阻等于 \SI{0.2}{\ohm}。与检流计并联一个电阻为 \SI{0.05}{\ohm} 的分流器（shunt）。应再与这一组合串联多大的电阻，才能使总电阻等于检流计本身的电阻？
\pnote{检流计即人教版所说的灵敏电流计。并联小电阻（分流器）以扩大量程、串联大电阻以改装电压表，正是两种电表的构造原理（人教版必修第三册第十一章第 4 节有讨论）。本题更进一步：串入一个适当的电阻，使总电阻恢复为检流计的原电阻——这样改装后的仪表接入电路时不改变原电路的电阻，是设计分流电路时常用的考虑。}

\item 一个电阻为 \SI{100}{\ohm} 的电压表，接在电动势为 \SI{10}{\volt}、内阻为 \SI{1}{\ohm} 的电池两端。求电压表的读数，并计算设其电阻为无穷大时读数的相对误差。

\item 一个电阻为 \SI{1}{\ohm} 的电流表，接入外电阻为 \SI{49}{\ohm}、电源电动势为 \SI{10}{\volt} 而内阻为 \SI{1}{\ohm} 的电路。求电流表的读数，并计算设其电阻为零时读数的相对误差。
\end{enumerate}

\cleardoublepage

% NEWTERMS: electromotive force (emf) | 电动势 | 正文全书首现于本章，首现括注（含简称 emf）；人教版必修第三册第十二章（记作 E）；ch-30 评注沿用，首现括注以本章为准
% NEWTERMS: seat of the electromotive force | 电动势源 | 底本旧式英语（即电源 источник ЭДС），首现括注；后文径用"电源"
% NEWTERMS: internal resistance | 内阻 | 首现括注；人教版亦简称内阻（记 r）；ch-30 沿用"内阻"
% NEWTERMS: portion of a circuit / element of a circuit | 电路的一段／一段电路 | 底本两词互用，首现分别括注；初中欧姆定律旧称"部分电路的欧姆定律"
% NEWTERMS: closed circuit | 闭合电路 | 首现括注；人教版必修第三册第十二章第 2 节《闭合电路的欧姆定律》
% NEWTERMS: voltage | 电压 | 正文首现括注；人教版记 U（原书 V）；初中已学
% NEWTERMS: potential drop | 电势降 | 首现括注；即"电压降"
% NEWTERMS: potential jump | 电势跃变 | 首现括注；发生在电池双电层内，机理见评注
% NEWTERMS: galvanic cell | 原电池 | 首现括注；与人教版化学同名术语一致
% NEWTERMS: external circuit | 外电路 | 首现括注；路端电压即外电路电压（"路端电压"定名在评注层）
% NEWTERMS: voltmeter | 电压表 | 首现括注；取人教版名（旧译伏特计）；理想化条件见正文
% NEWTERMS: ammeter | 电流表 | 习题首现括注；取人教版名（旧译安培计）
% NEWTERMS: galvanometer | 检流计 | 习题首现括注；名词委名，人教版称灵敏电流计
% NEWTERMS: shunt | 分流器 | 习题首现括注；电表改装部件
% NEWTERMS: current source | 电源 | 承 ch-30 定名（原书 источник тока 字面"电流源"，避与恒流源混淆）；本章习题 5 为正文首现，前经"电动势源（即电源）"已立名，不另括注
% NEWTERMS: resistor | 电阻 | 本章正文首现（R₁、R₂）；自明词不括注
% NEWTERMS: current | 电流 | 承 ch-26 章尾格言首现，不括注；人教版初中第十五章起
% NEWTERMS: relative error | 相对误差 | 常规数学术语，不括注
% NEWTERMS: approximation formula | 近似公式 | 自明词，不括注；(1+λ)^α ≈ 1+αλ 即二项式/泰勒一阶近似
% NEWTERMS: Ohm's law | 欧姆定律 | 承 ch-26 章尾格言首现；通行译名，不注
% NEWTERMS: terminal voltage | 路端电压 | 仅评注层引入（U=E−Ir）
% NEWTERMS: Ohm's law for a circuit portion containing an emf | 一段含源电路的欧姆定律 | 通行名，评注层给出；正文按底本作"适用于电路这一段的欧姆定律"
% NEWTERMS: equipotential points merging | 等势点合并（缩点） | 通行名，评注层给出；正文按底本作"电势相同的两点可合并成一点"
% NEWTERMS: dual layer (electric double layer) | 双电层 | 仅评注层引入
% NEWTERMS: Kirchhoff's second law | 基尔霍夫第二定律 | 仅评注层引入（回路电压方程）
% NEWTERMS: Problems (section title) | 习题 | 本章 5 题，接续全书连续编号（resume*=problems）
